\subsection{На каких дисциплинах основано}

Машинное обучение -- это наука, изучающая алгоритмы, автоматически улучшающиеся благодаря опыту.
Для построения таких методов используются средства математической статистики, численных методов,
математического анализа, методов оптимизации, теории вероятностей, теории графов, различные техники работы с
данными в цифровой форме.

\subsection{Разница между классическим программированием и ML}

Классическое программирование -- это автоматизация хода мыслей человека, а машинное обучение -- автоматизация
процесса вывода мышления из данных.
Более подробно в Таблице~\ref{tab:comp}

\begin{table}[h]
    \centering
    \begin{tabular}{>{\raggedright\arraybackslash}p{0.25\textwidth}
        | >{\raggedright\arraybackslash}p{0.33\textwidth} | >{\raggedright\arraybackslash}p{0.33\textwidth}}
        \toprule
        \textbf{Аспект} & \textbf{Классическое программирование} & \textbf{Машинное обучение} \\
        \midrule
        Основная задача & Выполнение четко определенной последовательности инструкций & Обобщение на основе данных
        и прогнозирование для новых, ранее невиданных случаев \\ \hline
        Правила & Явно задаются программистом в виде алгоритмов, условий (if-else), циклов & Выводится автоматически
        из данных в процессе обучения.
        Скрыта внутри параметров модели (например, весов нейронной сети) \\ \hline
        Входные данные & Часто структурированные (числа, строки, объекты) & Любые данные: таблицы, текст,
        изображения, звук.
        Качество и количество данных критически важны. \\ \hline
        Объект обучения & Нет.
        Программа статична, пока код не изменится & Параметры модели настраиваются в процессе обучения. \\ \hline
        Критерий успеха & Корректность выполнения алгоритма, отсутствие багов, производительность & Точность
        (accuracy), ошибка (loss) на новых данных (не тех, на которых учились). \\ \hline
        Гибкость & Низкая.
        Для новых условий или данных часто нужно переписывать код & Высокая.
        Модель может адаптироваться к изменениям в данных (до определенных пределов) \\ \hline
        Область применения & Детерминированные задачи с четкой логикой: веб-серверы, ОС, калькуляторы,
        бизнес-логика & Задачи с неявными правилами или большой изменчивостью: распознавание образов,
        прогнозирование, рекомендации, NLP, беспилотные автомобили \\ \hline
        Пример решения & Программист изучает кошек, пишет правила & Собираются сотни тысяч фото с кошками и
        без.
        Модель (например, CNN) обучается на них.
        Она сама находит значимые признаки и предсказывает вероятность, что на новом фото -- кошка. \\
        \bottomrule
    \end{tabular}
    \caption{Сравнительная таблица}
    \label{tab:comp}
\end{table}

\subsection{Базовые термины ML}

Типичная функция в машинке отображает \textbr{объекты или примеры (samples)} в \textbr{предсказания (targets)}.
Мы назовём функцию, отображающую объекты в предсказания, \textbr{моделью}, а имеющийся у нас набор примеров --
\textbr{обучающей выборкой или датасетом}.
Обучающая выборка состоит из: объектов и ответов, которые мы также будем иногда называть \textbr{таргетами}

Измеримые качества объекта обычно называются \textbr{признаками}, а совокупность свойств, которые мы выделили у
объекта – его \textbr{признаковым описанием}.
Вот несколько простых и распространённых разновидностей признаков:

\begin{itemize}

    \item Численные – например, рост или доход.
    Иногда отдельно выделяют вещественные и целочисленные признаки

    \item Категориальные признаки принимают значения из некоторого дискретного множества.
    Например, профессия человека или день недели

    \item Бинарные признаки принимают два значения: 0 и 1 или «да» и «нет».
    С ними можно работать и как с численными, и как с категориальными

    \item Среди категориальных признаков иногда выделяют ординальные.
    Они принимают значения из некоторого упорядоченного дискретного множества
    Класс опасности химического вещества (бывает от 1-го до 4-го) или год обучения студента являются ординальными

\end{itemize}

\textbr{feature engineering} -- процесс извлечения данных из признаков

\textbr{Матрица объекты-признаки} -- численная таблица, где строки которой соответствуют объектам, а столбцы --
признакам

\textbr{Предсказательные модели}, то есть модели вида $y=f(x)$, которые пытаются уловить зависимость между
признаковым описанием $x$ объекта и таргетом $y$.
Иногда мы будем иметь дело и с \textbr{моделями данных}: \(''\) признак имеет нормальное распределение \(''\).
Чаще всего предсказательные модели мы будем брать из некоторого параметрического семейства y$=f_w(x)$,
где $w$ — параметры, которые мы будем подбирать по данным

Кроме выбора модели важен также выбор алгоритма обучения.
\textbf{Алгоритм обучения} -- это процедура, которая превращает обучающую выборку в обученную модель.
Скажем, в примере выше для константной модели мы в качестве алгоритма обучения использовали поиск нуля градиента.
Градиентные методы используются для обучения многих моделей, и это очень богатый класс методов оптимизации, из
которого порой не так просто выбрать лучший

\textbf{Метрика} -- это способ численно измерить то, насколько модель хорошо справилась с задачей.
Метрика не единственная характеристика модели.
Ещё выделяют размер, скорость работы, интерпретируемость, отсутствие дискриминации данных